\section{Further criteria without Noetherian hypotheses}

\subsection{Bases with no proper injective epimorphism}

\begin{proposition}\label{prop:epimorphism-rigid-bases}
Every injective ring epimorphism $R\to C$ is an isomorphism in each
of the following situations:
\begin{enumerate}
\item $R$ is absolutely flat.
\item $R/N$ is absolutely flat and $N$ is a T-nilpotent ideal.
\item $R$ is a directed union of unital self-injective subrings.
\item $R$ is local with nil maximal ideal $\mathfrak m$, and there is
an ideal $I$ with $\mathfrak m^NI=0$ for some $N$, such that every
injective epimorphism from $R/I$ is an isomorphism.
\end{enumerate}
For any such $R$ and any inclusion $R\subset S$, unique polynomial
co-operations are equivalent to
\[
              (R:_S D_d(R))=R\quad\text{for every }d,
 \qquad\text{equivalently}\qquad\Omega(R,S)=R[X].
\]
Every positive case has all multivariable tensor isomorphisms.
\end{proposition}
\begin{proof}
In the absolutely flat case $C/R$ is flat, so tensoring
$0\to R\to C\to C/R\to0$ with any module preserves the initial
injection. Taking that module to be $C/R$, epicness gives
$C\otimes_R(C/R)=0$, hence $C/R=0$.

For (2), the nil ideal $NC$ contracts to $N$: if its contraction
were larger, its image in the reduced ring $R/N$ would contain a
nonzero nilpotent. The induced injective epimorphism
$R/N\to C/NC$ is therefore an isomorphism, giving $C=R+NC$.
For a fixed element of $C$, repeatedly expand each coefficient by
this equality. This produces a finitely branching tree; label its
edges by the elements of $N$ occurring in the expansions and prune
every branch whose product is zero. An infinite branch would
contradict T-nilpotence. K\"onig's lemma makes the remaining tree
finite, so the original element is a finite sum of elements of $R$.

For (3), write $R=\bigcup R_\alpha$. The equality
$c\otimes1=1\otimes c$ in $C\otimes_R C$ uses only finitely many
balancing relations, so it already holds over one $R_\alpha$.
Self-injectivity extends the identity of $R_\alpha$ to an
$R_\alpha$-linear retraction $\rho:C\to R_\alpha$. Applying
$\rho\otimes1$ gives $c=\rho(c)\in R$.

For (4), $\mathfrak mC$ is nil and therefore proper. Reduction gives
an epimorphism from the residue field to the nonzero algebra
$C/\mathfrak mC$, hence an isomorphism. Thus $C=R+\mathfrak mC$.
Multiplication by $I$ and finite iteration give
$IC=I+\mathfrak mIC=I+\mathfrak m^NIC=I$. Consequently
$R/I\subset C/I$ is an injective epimorphism, so $C/I=R/I$ and
$C=R$.

In every case apply the coefficient-ring obstruction to
$C=\bigcup_d(R:_S D_d(R))$. Tensor injectivity forces epicness,
and the preceding conclusions force $C=R$. The converse is the
ordinary polynomial case. The equivalence between absolute
flatness and reduced zero-dimensionality is standard
\cite[Tag~092F]{StacksAbsolutelyFlat}.
\end{proof}

\begin{corollary}\label{cor:boolean-and-unbounded-nilpotent-bases}
For a Boolean base ring $R$, unique polynomial co-operations for
$R\subset S$ occur exactly when $S=R$. The preceding proposition
also determines every coefficient extension of
\[
 R=\mathbb F_2[z,x_1,x_2,\ldots]/
       (z^2,zx_i,x_i^{2^{i+1}}:i\ge1).
\]
This ring is local, zero-dimensional and non-Noetherian. Its maximal
ideal is not T-nilpotent, and its elementwise nilpotence has no
uniform bound.
\end{corollary}
\begin{proof}
For a Boolean ring $D_2(R)=0$: every prime quotient is $\mathbb F_2$,
so every three-node determinant vanishes in all prime quotients,
and the ring is reduced. Thus the effective coefficient ring is
$S$, and the absolutely flat case applies.

For the displayed ring put
$A_n=\mathbb F_2[x_1,\ldots,x_n]/(x_i^{2^{i+1}})$ and $A=\bigcup A_n$.
The coefficient of the top monomial defines a perfect associative
pairing on $A_n$, pairing each monomial with its complement. It
identifies $A_n$ with its vector-space dual as an $A_n$-module.
That dual is injective because
$\operatorname{Hom}_{A_n}(M,\operatorname{Hom}_{\mathbb F_2}(A_n,\mathbb F_2))
\simeq\operatorname{Hom}_{\mathbb F_2}(M,\mathbb F_2)$ is exact.
Apply (3) to $A$, then (4) to $R=A\oplus\mathbb F_2z$ with
$I=\mathbb F_2z$ and $\mathfrak mI=0$.
Every element of $\mathfrak m$ uses finitely many variables and is
nilpotent. The independent classes of the $x_i$ in
$\mathfrak m/\mathfrak m^2$ show non-Noetherianity; the products
$x_1\cdots x_n$ never vanish, and $x_i$ has exact nilpotence order
$2^{i+1}$. These prove all remaining assertions.
\end{proof}

\Needspace{9\baselineskip}
\subsection{Regular interpolation and depth}

\begin{theorem}\label{thm:regular-vandermonde-criterion}
Let $Q(R)$ be the total quotient ring of a nonzero commutative ring
$R$. Suppose every $D_d(R)$ contains a regular element and
$(R:_{Q(R)}D_d(R))=R$. For every inclusion $R\subset S$, the
following are equivalent: unique polynomial co-operations; binary
tensor injectivity; faithful evaluation on $R$; the conditions
$\operatorname{Ann}_S D_d(R)=0$ for all $d$; and
$\Omega(R,S)=R[X]$. All positive cases have every tensor
isomorphism. The hypotheses hold for every infinite domain with
$\operatorname{Int}(R)=R[X]$, for every ring $A[Z,W]$, and for
$A[Z]$ whenever each $D_d(A)$ contains a regular element.
\end{theorem}
\begin{proof}
Localize $S$ at the regular elements of $R$. Cramer's rule says
$D_d(R)a_i\subset R$ for every coefficient of an integer-valued
polynomial. A regular element of this ideal makes its rationalized
coefficient lie in $Q(R)$, and the assumed colon identity puts it
in $R$. Rationalization therefore splits
\[
              \Omega(R,S)=R[X]\oplus\mathcal N,
\]
where $\mathcal N$ is the module of null polynomials. To identify
the kernel, clear the regular annihilators of the finitely many
coefficients; its values lie in $R$, so regularity forces them
to vanish. Conversely Cramer's rule kills each null coefficient
by $D_d(R)$ and therefore by a regular element.

If $\mathcal N\ne0$, choose a nonzero null coefficient $a$ with
$D_d(R)a=0$. The monic polynomial $P_d$ from the coefficient-ring
theorem has values in $D_d(R)$, so $aP_d$ is nonzero and null.
The tensor $aP_d\otimes P_d-P_d\otimes aP_d$ lies in the two
different mixed summands and is nonzero; extraction of the monic
coefficient detects its first component. Its polynomial image is
zero. This proves the injectivity implication. The other
equivalences follow from the faithful-evaluation criterion and
ordinary polynomials.

For an infinite domain, $D_d(R)\ne0$, and the coefficient theorem
identifies the colon condition with $\operatorname{Int}(R)=R[X]$.
For $R=A[Z]$, evaluate a rational integer-valued polynomial on
$A$-points. A regular $v\in D_d(A)$ bounds every coefficient in
$A[1/v][Z]$. Evaluate next at $Z^N$, choosing $N$ larger than all
coefficient degrees in $Z$. The degree intervals are disjoint;
the fact that the value belongs to $A[Z]$ forces all coefficients
to belong to $A[Z]$. Finally $D_d(A[W])$ contains a monic
polynomial: use the determinant on $0,1,W,\ldots,W^{d-1}$
(and $1$ for $d=1$). A monic polynomial is regular, proving
the assertion for arbitrary $A[Z,W]$.
\end{proof}

\begin{proposition}\label{prop:power-generated-and-regular-sequence-bases}
Let $R$ be local with finite residue field and maximal ideal
$\mathfrak m$.
If $\mathfrak m$ is generated by its $N$th powers for every $N$,
then unique polynomial co-operations are equivalent to
$(R:_S\mathfrak m)=R$, equivalently $\Omega(R,S)=R[X]$.
This applies in particular to local rings of characteristic $p$
with surjective Frobenius. If instead $\mathfrak m$ contains a
regular sequence of length two, the same equivalence holds with
$\Omega(R,S)=R[X]$, for every coefficient extension $S$.
\end{proposition}
\begin{proof}
In the power-generated case the determinant calculation in the
proof of Theorem~\ref{thm:valuation-conductor-arbitrary-coefficients}
uses only the local finite residue field and generation by powers:
$D_d(R)=R$ below the residue cardinality and $D_d(R)=\mathfrak m$
thereafter. The same proof with $C=(R:_S\mathfrak m)$ gives
$\mathfrak mC=\mathfrak m$ under epicness and then $C=R$.
Surjective Frobenius gives generation by every sufficiently large
$p$-power and consequently by every fixed $N$th power.

For a regular sequence $a,t$, both $a^e$ and $t^e$ belong to
$D_d(R)$ for a suitable $e$, by evaluating the same determinant
polynomial at $a$ and $t$. For any effective coefficient $b$,
write $x=a^eb\in R$ and $y=t^eb\in R$. Then $t^ex=a^ey$.
Regularity of $t^e$ modulo $a^e$ gives $x=a^ez$ for $z\in R$.
Rationalization at regular elements therefore sends $b$ to $z$.
This defines a ring retraction of the effective coefficient ring
to $R$. Under tensor injectivity that inclusion is epic, so the
retraction forces equality. The converse is ordinary polynomials.
\end{proof}

\subsection{A principal maximal ideal, valuation bases and binomial bases}

\begin{theorem}\label{thm:principal-maximal-coefficients}
Let $(R,\mathfrak m)$ be local with finite residue field and
$\mathfrak m=\pi R$. No regularity of $\pi$ is assumed. Put
$C=\{s\in S:\pi^ns\in R\text{ for some }n\}$.
If $\pi$ is nilpotent, unique polynomial co-operations occur
exactly when $S=R$. Otherwise they occur exactly when
\[
 C=R\quad\text{or}\quad
 \pi\text{ is regular on }R\text{ and }C=R[1/\pi]
 \text{ canonically inside }S.
\]
In every positive case all tensor maps are isomorphisms. If $\pi$
is regular on $R$, the criterion is equivalently
$\operatorname{Ann}_S(\pi)=0$.
\end{theorem}
\begin{proof}
If $\pi$ is nilpotent, the finite residue field and the filtration
by its powers make $R$ finite, so the finite-base theorem applies.
Assume otherwise and set $P=\bigcap_{n\ge0}\pi^nR$.
An element outside $P$ is a finite power of $\pi$ times a unit.
Consequently $\operatorname{Ann}_R(\pi)\subset P$, and
$P=\pi P$: if $x=\pi y\in P$, comparison with expressions
$x=\pi^{n+1}z$ gives $y\in\pi^nR+\operatorname{Ann}_R(\pi)$
for every $n$. It follows that $R/P$ is a DVR with parameter
$\bar\pi$. Indeed every nonzero element has a finite valuation,
products have the sum of valuations, and an ideal is generated
by an element of least valuation.

Finite grids of distinct powers of $\pi$ show that every coefficient
of an integer-valued polynomial lies in $C$. Conversely a power
of a monic polynomial null modulo $\pi R$ places every element of
$C$ in the effective coefficient ring. The ideals satisfy $PC=P$:
if $\pi^nc=a$ and $x\in P$, divide $x$ by arbitrarily high powers
of $\pi$ inside $P$ to obtain $xc\in P$. Reduction gives, in all
arities,
\[
 \Omega_n(R,C)/P\Omega_n(R,C)
       \simeq\Omega_n(R/P,C/P),\qquad
 P\Omega_n(R,C)=P[X_1,\ldots,X_n].
\]
Surjectivity follows by lifting coefficients; the preimage of
$R/P$ in $C$ is $R$. The kernel formula follows from $PC=P$.

For completeness the tensor kernel has the same elementary
description:
\[
 P\Omega(R,C)^{\otimes_Rn}\simeq P[X_1,\ldots,X_n].
\]
Given $h\in P$ times a pure tensor, write $h=\pi^Nh'$ with
$h'\in P$ and take $N$ large enough to clear the finitely many
coefficients of every factor, distributing this power among the
factors. Balancing reduces the expression to ordinary monomials
with coefficients in $P$. Conversely the polynomial coefficient
map on those ordinary monomials is injective. This proves the
displayed identification and shows that tensor injectivity and
surjectivity hold upstairs exactly when they hold after reduction.
Existence of either substituted tensor also lifts: correct any
lift by its ordinary polynomial error with coefficients in $P$.

The arbitrary-coefficient DVR case of the Noetherian theorem now
says $C/P=R/P$ or $C/P=\operatorname{Frac}(R/P)$ canonically.
The first case gives $C=R$. In the second, lift an inverse of
$\bar\pi$ to get $\pi c=1+p$, $p\in P$. Since $1+p$ is a unit
of $R$, $\pi$ is invertible in $C$, hence regular on the embedded
ring $R$, and $C=R[1/\pi]$. These prove necessity and sufficiency.
When $\pi$ is regular on $S$, $C$ embeds in $R[1/\pi]$; any
element outside $R$ has a numerator equal to a unit times a
finite power of $\pi$, so it generates $1/\pi$. Thus exactly
the two positive alternatives occur. Conversely either alternative
has no $\pi$-torsion, and every $\pi$-torsion element of $S$ lies
in $C$. This proves the last equivalence.
\end{proof}

\begin{corollary}\label{cor:valuation-arbitrary-coefficients}
Let $V$ be a valuation domain that is not a field. For any inclusion
$V\subset S$, unique polynomial co-operations occur exactly when
the residue field is infinite or
$\operatorname{Ann}_S(\mathfrak m_V)=0$. All positive cases have
every tensor isomorphism.
\end{corollary}
\begin{proof}
Infinite residue fields give ordinary polynomials. A principal
maximal ideal is covered by the preceding theorem, since its
generator is regular on $V$. If the maximal ideal is nonprincipal,
the ordered value group has no least positive element. Repeatedly
choosing a smaller positive value gives, for every $\epsilon>0$
and $N$, a $\delta>0$ with $N\delta<\epsilon$.
Thus $\mathfrak m_V$ is generated by its $N$th powers for every
$N$, and the power-generated criterion applies.
Here $(V:_S\mathfrak m_V)=V\oplus\operatorname{Ann}_S(\mathfrak m_V)$.
Indeed rationalization sends this colon into $K$: multiply by any
nonzero element of the maximal ideal. Its rational image has
nonnegative valuation, since arbitrarily small positive values
exclude a negative one. Subtract that element of $V$; the remainder
is killed by the whole maximal ideal. The sum is direct because
$V$ is a domain. This proves the equivalence.
\end{proof}

\begin{corollary}\label{cor:valuation-free-newton-basis}
Let $V$ be an infinite valuation domain with fraction field $K$,
and let $V\subset S$ be a torsion-free extension. Identify
$L=S\cap K$ inside $K\otimes_VS$. If the residue field is infinite,
or the maximal ideal is nonprincipal, then $\Omega(V,S)=V[X]$.
If the nonzero maximal ideal is $\pi V$ and the residue field has
$q$ elements, then
\[
 \Omega(V,S)=
 \begin{cases}
 V[X],&L=V,\\
 \operatorname{Int}(V),&L\ne V,
 \end{cases}
\]
and $\operatorname{Int}(V)$ has the free Newton basis described
below. Thus all these modules are free and every multivariable
polynomial tensor map is an isomorphism, in arbitrary valuation
rank.
\end{corollary}
\begin{proof}
Rationalization reduces the coefficient extension to $L$.
If $V$ is a field, or its residue field is infinite, Vandermonde
interpolation with pairwise unit differences gives $V[X]$.
Suppose its maximal ideal is nonprincipal. For every positive
value $\gamma$ and integer $N>0$ one can choose a positive value
$\delta$ with $N\delta<\gamma$. Indeed any positive value admits
a smaller one, and replacing $0<\xi<\epsilon$ by the smaller of
$\xi$ and $\epsilon-\xi$ gives a positive value at most
$\epsilon/2$ in the ordered-group sense. Finite iteration proves
the assertion without an Archimedean hypothesis.

If a coefficient of a degree-$d$ integer-valued polynomial had
value $-\gamma<0$, choose $a\in V$ with positive value $\delta$
such that $d(d+1)(d+2)\delta/6<\gamma$. The Vandermonde determinant
on $0,a,a^2,\ldots,a^d$ has exactly this smaller value: every
difference of positive powers has the smaller exponent's value.
Its adjugate bounds the negative of each coefficient value by
that determinant value, a contradiction. Hence
$\operatorname{Int}(V)=V[X]$.

In the remaining case choose representatives
$c_0=0,c_1,\ldots,c_{q-1}$ modulo $\pi V$ and, for
$n=\sum_jn_jq^j$, set $a_n=\sum_jc_{n_j}\pi^j$. Define
\[
 w_q(n)=\sum_{j\ge1}\lfloor n/q^j\rfloor,
 \qquad
 b_n(X)=\frac{\prod_{i<n}(X-a_i)}{\prod_{i<n}(a_n-a_i)}.
\]
First differing digits show that the denominator has value
$w_q(n)v(\pi)$. Every element of $V$ has an arbitrarily long
finite string of residue digits. At least $\lfloor n/q^j\rfloor$
of the first $n$ nodes share its first $j$ digits. Counting these
finitely many contributions shows that the numerator has at
least the denominator's value. This remains valid if a difference
has value greater than every integer multiple of $v(\pi)$.
Thus $b_n$ is integer-valued, vanishes at earlier nodes, and has
value $1$ at $a_n$. Successive evaluation gives
$\operatorname{Int}(V)=\bigoplus_{n\ge0}Vb_n$.

Each $b_n$ has coefficients in $V[1/\pi]$. If $b\in L\setminus V$,
the least positive value $v(\pi)$ satisfies $-v(b)\ge v(\pi)$,
so $(\pi b)^{-1}\in V$ and $\pi^{-1}=(\pi b)^{-1}b\in L$.
Consequently the entire basis is available over every proper
overring $L$. Iterating the triangular evaluation argument in
each variable gives the tensor product basis, proving all the
asserted tensor isomorphisms.
\end{proof}

\begin{corollary}\label{cor:binomial-arbitrary-coefficients}
Let $R$ be a nonzero binomial ring. For every inclusion $R\subset S$,
unique polynomial co-operations occur if and only if $S$ is
$\mathbb Z$-torsion-free. All positive cases have all multivariable
tensor isomorphisms. In particular this applies to
$\operatorname{Int}(\mathbb Z)$ and to
$\mathbb Z\oplus\mathbb Q\epsilon$, $\epsilon^2=0$.
\end{corollary}
\begin{proof}
Put $(X)_n=X(X-1)\cdots(X-n+1)$. The integral Stirling identities
give, without canceling integers in $S$,
\[
 F(r)=\sum_{n=0}^{\deg F}\Delta^nF(0)\binom rn,
 \qquad
 \Omega(R,S)=\bigoplus_{n\ge0}(R:_S n!)(X)_n.
\]
In several variables the coefficient ideal is
$(R:_S\prod_i n_i!)$. Preservation in a fixed degree is a finite
test on the integer grid, so these modules commute with arbitrary
localization. Localizations of $R$ remain binomial: degree induction
with $s^dF(X)-F(sX)$ proves preservation on localized inputs,
applied to $F=\binom Xn$ in $\mathbb Q\otimes R$.

At a local base of residue characteristic zero every integer is
a unit, giving ordinary polynomials. At residue characteristic
$p$, the identity $a^p-a\in pR$ makes the local ring $R/pR$
equal to $\mathbb F_p$: every element of its maximal ideal
satisfies $a(1-a^{p-1})=0$ with the second factor a unit.
Thus the maximal ideal is $pR$, and $p$ is regular on $R$.
The principal-maximal theorem gives precisely the condition
$\operatorname{Ann}_{S_{\mathfrak m}}(p)=0$.
These local conditions are equivalent to integer torsion-freeness
of $S$. A nonzero integer-torsion element supplies a nonzero
$p$-torsion element for some prime $p$; localize at a maximal ideal
containing its annihilator to detect it. Conversely regularity of
each integer persists under localization. The local tensor
isomorphisms glue by testing kernels and cokernels.

The ring $\operatorname{Int}(\mathbb Z)$ is binomial because
$\binom{f(n)}j$ is an integer for every integer $n$. For
$r=a+b\epsilon\in\mathbb Z\oplus\mathbb Q\epsilon$,
$\binom rj=\binom aj+b(\binom Xj)'(a)\epsilon$ lies in the same
ring. Its square-zero ideal is not finitely generated over the
base, so this example is non-Noetherian as well as nonreduced.
\end{proof}

\subsection{A bounded nilpotent maximal ideal}

\begin{theorem}\label{thm:bounded-p-nilpotence}
Let $R$ be a local $\mathbb F_p$-algebra with finite residue field
$k=\mathbb F_q$, and suppose $x^p=0$ for every $x\in\mathfrak m$.
For any inclusion $R\subset S$, unique polynomial co-operations
are equivalent to the two conditions
\[
                 \mu:S\otimes_R S\longrightarrow S
                       \text{ is an isomorphism},
                 \qquad(R:_S\mathfrak m)=R.
\]
Here $\mu(s\otimes t)=st$ is the multiplication map.
All positive cases have every tensor isomorphism. If
$\mathfrak m$ is nilpotent, these conditions force $S=R$.
\end{theorem}
\begin{proof}
Frobenius identifies its image $R^p$ with the residue field,
giving a coefficient field. Put $Q=(X^q-X)^p$, a null polynomial
on $R$, and choose the Hermite basis $h_{a,j}$, $a\in k$, $0\le j<p$,
dual to the Hasse derivatives of order below $p$ at these points.
Write $C_j=(R:_S\mathfrak m^j)$. Monic division and Hermite
interpolation give
\[
                     \Omega(R,S)=\bigoplus_{a,j}C_jh_{a,j}
                                      \oplus Q S[X].
\]
To check the value condition, evaluate at $a+\lambda n$ for
$\lambda\in\mathbb F_p$ and $n\in\mathfrak m$. The unit
Vandermonde matrix separates the coefficients of $n^j$.
Polarization, using that $j!$ is a unit for $j<p$, shows that
these powers generate $\mathfrak m^j$. This proves the formula.

Every element of $S$ is a leading coefficient of $cQ$, so tensor
injectivity makes $R\to S$ epic. For $c\in C_1$ and $x\in\mathfrak m$,
$(xc)^p=0$, whence $xc\in\mathfrak m$. Thus $\mathfrak mC_1=\mathfrak m$
and $C_1$ is a ring. The split Hermite summand $C_1h_{a,1}$ makes
$C_1\otimes_R C_1\to C_1$ injective under tensor injectivity.
Reduction of this epimorphism modulo $\mathfrak m$ gives the
identity residue field, forcing $C_1=R$.

Conversely $C_1=R$ implies all $C_j=R$ by iterated colons.
In several variables, mixed differences, the same scalar
interpolation and polarization put every low Hermite coefficient
in $(R:_S\mathfrak m^{\sum j_i})=R$. A term divisible by $Q$ in
at least one variable has arbitrary coefficient in $S$. The tensor
power of the displayed splitting has exactly these coefficient
modules because $S\otimes_R S=S$, proving every tensor
isomorphism. If $\mathfrak m^N=0$, the iterated colon identity
gives $S=(R:_S\mathfrak m^N)=R$.
\end{proof}
